\section{Practical Session 1: Phasor-Domain Analysis}

\subsection{Exercises}
\begin{enumerate}
    \item \textbf{Phasor Representation of Sinusoidal Voltages}\footnote{Exercises 2.1, 2.2, 2.5 and 2.9 from Ned Mohan's book \textit{"Electric power systems, a first course"}.}
    
    Express the following voltages as phasors:
    \begin{enumerate}
        \item $v_1(t) = \sqrt{2} \times 100 \cos(\omega t - 30^\circ)\ \text{V}$
        \item $v_2(t) = \sqrt{2} \times 100 \cos(\omega t + 30^\circ)\ \text{V}$
    \end{enumerate}

    \item \textbf{Series $R$-$L$-$C$ Circuit}
    
    The series $R$-$L$-$C$ circuit shown in Figure~\ref{fig:tp1_rlc} is in sinusoidal steady state at a frequency of $f = 60\ \text{Hz}$. The parameters are:
    \[
    V = 120\ \text{V},\quad R = 1.5\ \Omega,\quad L = 20\ \text{mH},\quad C = 100\ \mu\text{F}.
    \]
    Calculate the current $i(t)$ in this circuit by using phasor-domain analysis.

    \begin{figure}[htbp]
        \centering
        \begin{circuitikz}[american, scale=1.1]
            \draw (0,0) to[sV, v_>=$v(t){=}\sqrt{2}V\cos(\omega t)$] (0,3)
                  to[short, i=$i(t)$] (2,3)
                  to[L, l=$L$] (3.5,3)
                  to[R, l=$R$] (3.5,1.5)
                  to[C, l=$C$] (3.5,0)
                  -- (0,0);
        \end{circuitikz}
        \caption{Series $R$-$L$-$C$ circuit.}
        \label{fig:tp1_rlc}
    \end{figure}

    \item \textbf{Series-Parallel $R$-$L$-$C$ Circuit and Reactive Power Conservation}
    
    To the circuit of Figure~\ref{fig:tp1_ex4}, if a voltage of $100\angle 0^\circ\ \text{V}$ is applied, calculate the active power $P$, the reactive power $Q$, and the power factor $\cos\phi$. Show that:
    \[
    Q = \sum_k I_k^2 X_k
    \]

    \begin{figure}[htbp]
        \centering
        \begin{circuitikz}[american, scale=1.1]
            \draw (0,0) to[sV, v_>=$\bar{V}{=}100\angle 0^\circ\ \text{V}$] (0,2.5)
                  to[L, l=$j0.1\ \Omega$, i=$\bar{I}_L$] (3,2.5)
                  -- (4.5,2.5);
            \draw (3,2.5) to[C, l=$-j5\ \Omega$, i=$\bar{I}_C$] (3,0);
            \draw (4.5,2.5) to[R, l=$2\ \Omega$, i=$\bar{I}_R$] (4.5,0);
            \draw (0,0) -- (4.5,0);
        \end{circuitikz}
        \caption{Series-parallel circuit for Exercise 3.}
        \label{fig:tp1_ex4}
    \end{figure}

    \item \textbf{Power Factor Correction}
    
    In the circuit of Figure~\ref{fig:tp1_pfc}, the complex power drawn by the load impedance is:
    \[
    S_L = P_L + j Q_L = (1858.4 + j1031.3)\ \text{VA}.
    \]
    Calculate the capacitive reactance $X_C$ connected in parallel necessary to make the overall power factor equal to $0.9$ (leading), assuming the applied voltage has an RMS value of $V = 120\ \text{V}$ at $f = 60\ \text{Hz}$. Determine the required capacitance $C$.

    \begin{figure}[htbp]
        \centering
        \begin{circuitikz}[american, scale=1.1]
            \draw (0,0) to[sV, v_>=$\bar{V}_1$] (0,3)
                  to[short, i=$P{=}P_L$] (2.5,3)
                  -- (4.5,3);
            \draw (2.5,3) to[C, l=$-j Q_C$, i>^=$Q_C$] (2.5,0);
            \draw (4.5,3) to[generic, l={$P_L + j Q_L$}, i>^={$P_L + j Q_L$}] (4.5,0);
            \draw (0,0) -- (4.5,0);
        \end{circuitikz}
        \caption{Power factor correction circuit.}
        \label{fig:tp1_pfc}
    \end{figure}
\end{enumerate}

\newpage
\subsection{Solutions}

\subsubsection*{Solution to Exercise 1: Phasor Representation}
Using the cosine reference with RMS value as phasor magnitude:
\begin{align*}
    v(t) = \sqrt{2} V_{\text{RMS}} \cos(\omega t + \theta) \iff \bar{V} = V_{\text{RMS}} e^{j\theta} = V_{\text{RMS}} \angle \theta
\end{align*}

\begin{enumerate}[(a)]
    \item For $v_1(t) = \sqrt{2} \times 100 \cos(\omega t - 30^\circ)\ \text{V}$:
    \begin{equation}
        \boxed{\bar{V}_1 = V_{\text{RMS}} e^{j\theta} = \frac{V_m}{\sqrt{2}} e^{-j 30^\circ} = 100\ \text{V} \angle -30^\circ}
    \end{equation}

    \item For $v_2(t) = \sqrt{2} \times 100 \cos(\omega t + 30^\circ)\ \text{V}$:
    \begin{equation}
        \boxed{\bar{V}_2 = V_{\text{RMS}} e^{j\theta} = \frac{V_m}{\sqrt{2}} e^{+j 30^\circ} = 100\ \text{V} \angle +30^\circ}
    \end{equation}
\end{enumerate}

\subsubsection*{Solution to Exercise 2: Series $R$-$L$-$C$ Circuit}
\paragraph{Data:}
\begin{itemize}
    \item Voltage: $\bar{V} = 120\ \text{V}\angle 0^\circ$
    \item Resistance: $R = 1.5\ \Omega$
    \item Inductance: $L = 20\ \text{mH} = 20 \times 10^{-3}\ \text{H}$
    \item Capacitance: $C = 100\ \mu\text{F} = 100 \times 10^{-6}\ \text{F}$
    \item Frequency: $f = 60\ \text{Hz} \implies \omega = 2\pi f = 2\pi \times 60 = 376.99\ \text{rad/s}$
\end{itemize}

\paragraph{Impedance $Z$:}
The total impedance of the series circuit is:
\begin{align}
    Z &= R + j\omega L + \frac{1}{j\omega C} = R + j\left(\omega L - \frac{1}{\omega C}\right) \nonumber\\
      &= 1.5 + j\left(376.99 \times 20 \times 10^{-3} - \frac{1}{376.99 \times 100 \times 10^{-6}}\right) \nonumber\\
      &= 1.5 + j(7.5398 - 26.5258) = 1.5 - j18.986\ \Omega \nonumber\\
      &= 19.049\ \Omega \angle -85.48^\circ\quad (-1.492\ \text{rad})
\end{align}

\paragraph{Current Phasor $\bar{I}$:}
By Ohm's law:
\begin{equation}
    \bar{I} = \frac{\bar{V}}{Z} = \frac{120\angle 0^\circ}{19.049\angle -85.48^\circ} = 6.30\ \text{A} \angle +85.48^\circ\quad (+1.492\ \text{rad})
\end{equation}

\paragraph{Time-Domain Current $i(t)$:}
Converting back from the RMS phasor to the time domain:
\begin{equation}
    \boxed{i(t) = 6.30\sqrt{2} \cos(376.99 t + 1.492)\ \text{A} = 8.91 \cos(376.99 t + 85.5^\circ)\ \text{A}}
\end{equation}

\subsubsection*{Solution to Exercise 3: Series-Parallel Circuit and Power Conservation}
\paragraph{Data:}
\begin{itemize}
    \item Applied voltage: $\bar{V} = 100\ \text{V} \angle 0^\circ$
    \item Inductor impedance: $Z_L = j0.1\ \Omega$
    \item Capacitor impedance: $Z_C = -j5\ \Omega$
    \item Resistor impedance: $R = 2\ \Omega$
\end{itemize}

\paragraph{Total Equivalent Impedance $Z$:}
The capacitor and resistor are connected in parallel:
\begin{align}
    Z_p = Z_C \parallel R = \frac{R \cdot Z_C}{R + Z_C} = \frac{2(-j5)}{2 - j5} = \frac{-j10(2 + j5)}{2^2 + (-5)^2} = \frac{50 - j20}{29} = 1.7241 - j0.6897\ \Omega
\end{align}
Adding the series inductor:
\begin{align}
    Z &= Z_L + Z_p = j0.1 + (1.7241 - j0.6897) = 1.7241 - j0.5897\ \Omega \nonumber\\
      &= 1.822\ \Omega \angle -18.88^\circ
\end{align}

\paragraph{Complex Power $S$:}
\begin{align}
    S &= \bar{V} \bar{I}^* = \bar{V} \left(\frac{\bar{V}}{Z}\right)^* = \frac{V^2}{Z^*} = \frac{100^2}{1.822\angle +18.88^\circ} \nonumber\\
      &= 5487.93\ \text{VA} \angle -18.88^\circ
\end{align}
Since $S = P + jQ$:
\begin{align}
    \boxed{P = \text{Re}\{S\} = 5487.93 \cos(-18.88^\circ) = 5192.54\ \text{W}} \\
    \boxed{Q = \text{Im}\{S\} = 5487.93 \sin(-18.88^\circ) = -1775.89\ \text{var}}
\end{align}

\paragraph{Power Factor $\cos\phi$:}
\begin{equation}
    \boxed{\cos\phi = \frac{P}{|S|} = \frac{5192.54}{5487.93} = 0.946\quad (\text{leading})}
\end{equation}
Because $Q < 0$ (and $\phi < 0$), the circuit is capacitive overall: it produces net reactive power, meaning that the current phasor leads the voltage phasor.

\paragraph{Proof that $Q = \sum_k I_k^2 X_k$:}
We compute the branch currents:
\begin{itemize}
    \item Current through the series inductor:
    \[
    \bar{I}_L = \bar{I} = \frac{\bar{V}}{Z} = \frac{100\angle 0^\circ}{1.822\angle -18.88^\circ} = 54.88\ \text{A} \angle +18.88^\circ \implies I_L = 54.88\ \text{A}
    \]
    \item Current through the capacitor (using current divider):
    \begin{align*}
        \bar{I}_C &= \frac{R}{R + Z_C} \bar{I}_L = \frac{2}{2 - j5} (54.88\angle 18.88^\circ) \\
                  &= \frac{2}{5.385\angle -68.20^\circ} (54.88\angle 18.88^\circ) = 20.38\ \text{A} \angle +87.08^\circ \implies I_C = 20.38\ \text{A}
    \end{align*}
\end{itemize}

Calculating the individual reactive powers:
\begin{itemize}
    \item Resistor: $Q_R = 0\ \text{var}$
    \item Capacitor: $Q_C = X_C I_C^2 = (-5) \times (20.38)^2 = -2077.06\ \text{var}$
    \item Inductor: $Q_L = X_L I_L^2 = 0.1 \times (54.88)^2 = +301.17\ \text{var}$
\end{itemize}
Summing the individual contributions:
\begin{equation}
    \sum_k I_k^2 X_k = Q_R + Q_C + Q_L = 0 - 2077.06 + 301.17 = -1775.89\ \text{var} = Q
\end{equation}
This confirms that $Q = \sum_k I_k^2 X_k$. \hfill $\blacksquare$

\subsubsection*{Solution to Exercise 4: Power Factor Correction}
\paragraph{Data:}
\begin{itemize}
    \item Load complex power: $S_L = P_L + j Q_L = (1858.4 + j1031.3)\ \text{VA}$
    \item Applied voltage: $V = 120\ \text{V}$ (RMS)
    \item Desired overall power factor: $\cos\phi = 0.9$ (leading)
    \item Grid frequency: $f = 60\ \text{Hz} \implies \omega = 2\pi f = 376.99\ \text{rad/s}$
\end{itemize}

\paragraph{Determination of Total Reactive Power $Q_{\text{tot}}$:}
Connecting a pure capacitor in parallel does not alter the active power consumed by the system:
\[
P_{\text{tot}} = P_L = 1858.4\ \text{W}
\]
The new total complex power is $S = P_{\text{tot}} + j Q_{\text{tot}}$ with $Q_{\text{tot}} = Q_L + Q_C$.
Using the power factor relationship:
\begin{align}
    \cos\phi = \frac{P_{\text{tot}}}{|S|} = \frac{P}{\sqrt{P^2 + Q_{\text{tot}}^2}} \iff Q_{\text{tot}}^2 = P^2 \left( \frac{1}{\cos^2\phi} - 1 \right)
\end{align}
Taking the square root:
\begin{equation}
    Q_{\text{tot}} = \pm P \sqrt{\frac{1}{\cos^2\phi} - 1} = \pm 1858.4 \sqrt{\frac{1}{0.9^2} - 1} = \pm 900.064\ \text{var}
\end{equation}
Because the corrected power factor is required to be \textbf{leading}, the total reactive power must be negative:
\begin{equation}
    Q_{\text{tot}} = -900.064\ \text{var}
\end{equation}

\paragraph{Capacitor Reactive Power $Q_C$:}
Since $Q_{\text{tot}} = Q_L + Q_C$:
\begin{equation}
    Q_C = Q_{\text{tot}} - Q_L = -900.064 - 1031.3 = -1931.364\ \text{var}
\end{equation}

\paragraph{Capacitive Reactance $X_C$:}
Across a voltage $V = 120\ \text{V}$, the reactive power is related to the reactance by $Q_C = \frac{V^2}{X_C}$:
\begin{equation}
    \boxed{X_C = \frac{V^2}{Q_C} = \frac{120^2}{-1931.364} = -7.456\ \Omega \approx -7.46\ \Omega}
\end{equation}

\paragraph{Capacitance $C$:}
Since $X_C = -\frac{1}{\omega C}$:
\begin{equation}
    \boxed{C = \frac{1}{\omega |X_C|} = \frac{1}{2\pi \times 60 \times 7.456} = 3.557 \times 10^{-4}\ \text{F} \approx 0.356\ \text{mF} = 356\ \mu\text{F}}
\end{equation}
